\documentclass[10pt,oneside]{report}
\usepackage{soc2style}
\begin{document}
\doctitle{Criterion by Criterion Study Guide}{Sixty-one evidence-centered study cards and scenario exercises}{SOC2-02}
\section*{Operating convention and boundaries}
An AICPA criterion identifies an outcome against which controls may be evaluated. The selected point-of-focus headings identify relevant characteristics; they are not discrete certification checklist items. Every owner assignment, technological mechanism and audit test in this suite is a proposed example to be tailored to the system description, commitments, risk profile, period, and service organization responsibilities.

\begin{itemize}
\item Always establish system boundaries, the in-scope trust services categories, service commitments and system requirements first.
\item For a security-only engagement, use the common criteria. For additional trust services categories, include the corresponding additional criteria, subject to properly justified applicability.
\item Document complementary user entity controls and subservice organizations; do not treat vendor assurance reports as replacements for management ownership.
\item Privacy points of focus marked [P] or [C] describe typical processor or controller relevance; assess actual responsibilities rather than treating the annotations as absolute.
\end{itemize}

\chapter{How to reason about a criterion}
Treat each criterion as an assurance question: What objective is at risk? Which designed controls address that risk? Which of those controls operated during the relevant period? What complete and reliable evidence supports the conclusion? Avoid mistaking point-of-focus text for an exhaustive mandatory checklist.

\subsection{Study card technique}
\begin{enumerate}
\item Restate the published criterion in your own words without deleting its outcome or engagement qualifiers.
\item Describe a realistic failure mode that would violate a customer commitment.
\item Sketch one preventive and one detective implementation option, without declaring either mandatory.
\item Choose a population and a reproducible, time-bound evidence source.
\item State a test and the conditions under which an exception would require analysis.
\end{enumerate}

\chapter{Criterion-by-criterion study cards}
\section{Control Environment}
\Needspace{6\baselineskip}
\subsection{CC1.1 -- Ethics and conduct}
\begin{aicpabox}
\textbf{CC1.1}. COSO Principle 1: The entity demonstrates a commitment to integrity and ethical values.\par\sourcepage{14}
\end{aicpabox}
Engineering interpretation (example): Publish and annually attest a code of conduct; operate an independent reporting and escalation path.

Evidence reasoning: Approved conduct policy; attestations; case register with redactions. Use this only after confirming the control is part of the actual system.

Assessment exercise: Select hires and incumbent population; inspect communication, acknowledgment and deviations. Identify the population, sample unit, expected outcome and exception-handling process.

Common misconception: An unsigned policy without communicated expectations does not show operation.

Reflection prompt: If the control were disabled for two weeks, which indicators, owners, records and customer commitments would reveal the failure?

\Needspace{6\baselineskip}
\subsection{CC1.2 -- Board oversight}
\begin{aicpabox}
\textbf{CC1.2}. COSO Principle 2: The board of directors demonstrates independence from management and exercises oversight of the development and performance of internal control.\par\sourcepage{15}
\end{aicpabox}
Engineering interpretation (example): Schedule risk and control oversight sessions with documented challenges, decisions and assigned actions.

Evidence reasoning: Board or committee minutes; risk packs; action tracker. Use this only after confirming the control is part of the actual system.

Assessment exercise: Trace material control exceptions to board reporting and closure. Identify the population, sample unit, expected outcome and exception-handling process.

Common misconception: Management-only approval is not necessarily independent oversight.

Reflection prompt: If the control were disabled for two weeks, which indicators, owners, records and customer commitments would reveal the failure?

\Needspace{6\baselineskip}
\subsection{CC1.3 -- Structures and responsibility}
\begin{aicpabox}
\textbf{CC1.3}. COSO Principle 3: Management establishes, with board oversight, structures, reporting lines, and appropriate authorities and responsibilities in the pursuit of objectives.\par\sourcepage{16}
\end{aicpabox}
Engineering interpretation (example): Maintain organization chart, RACI, control ownership and vendor oversight responsibility.

Evidence reasoning: Org chart; role descriptions; RACI; delegations. Use this only after confirming the control is part of the actual system.

Assessment exercise: Sample high-risk controls for assigned accountable owner and segregation. Identify the population, sample unit, expected outcome and exception-handling process.

Common misconception: Unowned inherited cloud services create gaps across boundaries.

Reflection prompt: If the control were disabled for two weeks, which indicators, owners, records and customer commitments would reveal the failure?

\Needspace{6\baselineskip}
\subsection{CC1.4 -- Competence and training}
\begin{aicpabox}
\textbf{CC1.4}. COSO Principle 4: The entity demonstrates a commitment to attract, develop, and retain competent individuals in alignment with objectives.\par\sourcepage{17}
\end{aicpabox}
Engineering interpretation (example): Use role profiles, background screening where lawful, onboarding, continuing training and succession plans.

Evidence reasoning: Role matrix; training records; screening attestations. Use this only after confirming the control is part of the actual system.

Assessment exercise: Sample privileged roles for qualification and current training. Identify the population, sample unit, expected outcome and exception-handling process.

Common misconception: Completion counts alone do not establish role-appropriate capability.

Reflection prompt: If the control were disabled for two weeks, which indicators, owners, records and customer commitments would reveal the failure?

\Needspace{6\baselineskip}
\subsection{CC1.5 -- Accountability and performance}
\begin{aicpabox}
\textbf{CC1.5}. COSO Principle 5: The entity holds individuals accountable for their internal control responsibilities in the pursuit of objectives.\par\sourcepage{18}
\end{aicpabox}
Engineering interpretation (example): Tie control responsibilities to performance review, corrective action and incentives with safe escalation.

Evidence reasoning: Review rubric; acknowledgments; corrective-action logs. Use this only after confirming the control is part of the actual system.

Assessment exercise: Trace missed control duties to accountability and follow-up. Identify the population, sample unit, expected outcome and exception-handling process.

Common misconception: Do not use incentives that suppress incident reporting.

Reflection prompt: If the control were disabled for two weeks, which indicators, owners, records and customer commitments would reveal the failure?

\section{Information and Communication}
\Needspace{6\baselineskip}
\subsection{CC2.1 -- Quality information and inventory}
\begin{aicpabox}
\textbf{CC2.1}. COSO Principle 13: The entity obtains or generates and uses relevant, quality information to support the functioning of internal control.\par\sourcepage{19}
\end{aicpabox}
Engineering interpretation (example): Inventory services, repositories, assets, vendors and data flows with ownership, classification and periodic reconciliation.

Evidence reasoning: CMDB; data-flow diagram; classification register; reconciliation logs. Use this only after confirming the control is part of the actual system.

Assessment exercise: Reconcile sampled production resources against inventory and classifications. Identify the population, sample unit, expected outcome and exception-handling process.

Common misconception: Unknown cloud resources and shadow repositories make evidence incomplete.

Reflection prompt: If the control were disabled for two weeks, which indicators, owners, records and customer commitments would reveal the failure?

\Needspace{6\baselineskip}
\subsection{CC2.2 -- Internal communications}
\begin{aicpabox}
\textbf{CC2.2}. COSO Principle 14: The entity internally communicates information, including objectives and responsibilities for internal control, necessary to support the functioning of internal control.\par\sourcepage{21}
\end{aicpabox}
Engineering interpretation (example): Publish control objectives, incident reporting channels, policy changes and security awareness to relevant personnel.

Evidence reasoning: Communications log; trainings; incident channel instructions. Use this only after confirming the control is part of the actual system.

Assessment exercise: Sample changes and staff for timely notification and comprehension. Identify the population, sample unit, expected outcome and exception-handling process.

Common misconception: A training slide without attendance or targeting is weak evidence.

Reflection prompt: If the control were disabled for two weeks, which indicators, owners, records and customer commitments would reveal the failure?

\Needspace{6\baselineskip}
\subsection{CC2.3 -- External communications}
\begin{aicpabox}
\textbf{CC2.3}. COSO Principle 15: The entity communicates with external parties regarding matters affecting the functioning of internal control.\par\sourcepage{22}
\end{aicpabox}
Engineering interpretation (example): Define accurate customer, regulator, vendor and incident contact channels and commitments.

Evidence reasoning: Service agreements; trust communications; contact playbook. Use this only after confirming the control is part of the actual system.

Assessment exercise: Select applicable customer commitments and inspect external communication. Identify the population, sample unit, expected outcome and exception-handling process.

Common misconception: Misleading public descriptions can conflict with implemented controls.

Reflection prompt: If the control were disabled for two weeks, which indicators, owners, records and customer commitments would reveal the failure?

\section{Risk Assessment}
\Needspace{6\baselineskip}
\subsection{CC3.1 -- Objectives and risk appetite}
\begin{aicpabox}
\textbf{CC3.1}. COSO Principle 6: The entity specifies objectives with sufficient clarity to enable the identification and assessment of risks relating to objectives.\par\sourcepage{24}
\end{aicpabox}
Engineering interpretation (example): Record system objectives, commitments, risk tolerances and system boundaries before selecting safeguards.

Evidence reasoning: Risk appetite; system description; SLA register. Use this only after confirming the control is part of the actual system.

Assessment exercise: Trace each material commitment into risk identification and controls. Identify the population, sample unit, expected outcome and exception-handling process.

Common misconception: Generic security goals cannot prove customer-specific commitments.

Reflection prompt: If the control were disabled for two weeks, which indicators, owners, records and customer commitments would reveal the failure?

\Needspace{6\baselineskip}
\subsection{CC3.2 -- Risk identification and analysis}
\begin{aicpabox}
\textbf{CC3.2}. COSO Principle 7: The entity identifies risks to the achievement of its objectives across the entity and analyzes risks as a basis for determining how the risks should be managed.\par\sourcepage{27}
\end{aicpabox}
Engineering interpretation (example): Maintain risk inventory, threat scenarios, likelihood and impact, treatment and periodic re-evaluation.

Evidence reasoning: Risk register; workshop results; treatment approvals. Use this only after confirming the control is part of the actual system.

Assessment exercise: Sample assets/services and verify risk analysis and owners. Identify the population, sample unit, expected outcome and exception-handling process.

Common misconception: Risk scoring without linked response or reassessment is incomplete.

Reflection prompt: If the control were disabled for two weeks, which indicators, owners, records and customer commitments would reveal the failure?

\Needspace{6\baselineskip}
\subsection{CC3.3 -- Fraud risk}
\begin{aicpabox}
\textbf{CC3.3}. COSO Principle 8: The entity considers the potential for fraud in assessing risks to the achievement of objectives.\par\sourcepage{28}
\end{aicpabox}
Engineering interpretation (example): Model insider misuse, collusion, credential abuse and financial misrepresentation; design anti-fraud controls.

Evidence reasoning: Fraud scenarios; SoD review; whistleblower handling. Use this only after confirming the control is part of the actual system.

Assessment exercise: Trace identified fraud scenarios to mitigations and monitoring. Identify the population, sample unit, expected outcome and exception-handling process.

Common misconception: Do not equate fraud assessment with only malware detection.

Reflection prompt: If the control were disabled for two weeks, which indicators, owners, records and customer commitments would reveal the failure?

\Needspace{6\baselineskip}
\subsection{CC3.4 -- Significant change assessment}
\begin{aicpabox}
\textbf{CC3.4}. COSO Principle 9: The entity identifies and assesses changes that could significantly impact the system of internal control.\par\sourcepage{29}
\end{aicpabox}
Engineering interpretation (example): Trigger risk reviews for material architecture, acquisition, regulation and vendor changes.

Evidence reasoning: Change intake; impact assessments; risk decisions. Use this only after confirming the control is part of the actual system.

Assessment exercise: Sample major changes and assess whether risk review preceded deployment. Identify the population, sample unit, expected outcome and exception-handling process.

Common misconception: Operational changes may invalidate previously sufficient controls.

Reflection prompt: If the control were disabled for two weeks, which indicators, owners, records and customer commitments would reveal the failure?

\section{Monitoring Activities}
\Needspace{6\baselineskip}
\subsection{CC4.1 -- Control monitoring}
\begin{aicpabox}
\textbf{CC4.1}. COSO Principle 16: The entity selects, develops, and performs ongoing and/or separate evaluations to ascertain whether the components of internal control are present and functioning.\par\sourcepage{30}
\end{aicpabox}
Engineering interpretation (example): Maintain a monitoring plan mixing continuous telemetry and independent targeted reviews.

Evidence reasoning: Control monitoring schedule; dashboards; findings. Use this only after confirming the control is part of the actual system.

Assessment exercise: Compare monitoring coverage and executed reviews to control inventory. Identify the population, sample unit, expected outcome and exception-handling process.

Common misconception: Continuous alerts without ownership and review are not evaluations.

Reflection prompt: If the control were disabled for two weeks, which indicators, owners, records and customer commitments would reveal the failure?

\Needspace{6\baselineskip}
\subsection{CC4.2 -- Deficiency communication}
\begin{aicpabox}
\textbf{CC4.2}. COSO Principle 17: The entity evaluates and communicates internal control deficiencies in a timely manner to those parties responsible for taking corrective action, including senior management and the board of directors, as appropriate.\par\sourcepage{31}
\end{aicpabox}
Engineering interpretation (example): Grade deficiencies, route them to accountable owners, track deadlines and validate closure.

Evidence reasoning: Findings log; escalation tickets; validation artifacts. Use this only after confirming the control is part of the actual system.

Assessment exercise: Sample overdue issues for escalation and independent closure verification. Identify the population, sample unit, expected outcome and exception-handling process.

Common misconception: Closing a ticket without confirming risk mitigation is insufficient.

Reflection prompt: If the control were disabled for two weeks, which indicators, owners, records and customer commitments would reveal the failure?

\section{Control Activities}
\Needspace{6\baselineskip}
\subsection{CC5.1 -- Risk-based control activities}
\begin{aicpabox}
\textbf{CC5.1}. COSO Principle 10: The entity selects and develops control activities that contribute to the mitigation of risks to the achievement of objectives to acceptable levels.\par\sourcepage{31}
\end{aicpabox}
Engineering interpretation (example): Choose preventive/detective and manual/automated controls tied to documented risks.

Evidence reasoning: Risk-control matrix; approvals; control narratives. Use this only after confirming the control is part of the actual system.

Assessment exercise: Walk through sampled risks and inspect implemented control activity. Identify the population, sample unit, expected outcome and exception-handling process.

Common misconception: Checklist compliance does not replace risk-based design.

Reflection prompt: If the control were disabled for two weeks, which indicators, owners, records and customer commitments would reveal the failure?

\Needspace{6\baselineskip}
\subsection{CC5.2 -- Technology general controls}
\begin{aicpabox}
\textbf{CC5.2}. COSO Principle 11: The entity also selects and develops general control activities over technology to support the achievement of objectives.\par\sourcepage{32}
\end{aicpabox}
Engineering interpretation (example): Establish technology operations, infrastructure, access, development and maintenance governance.

Evidence reasoning: ITGC inventory; platform standards; administrative review. Use this only after confirming the control is part of the actual system.

Assessment exercise: Test technology dependencies of critical automated controls. Identify the population, sample unit, expected outcome and exception-handling process.

Common misconception: Strong application controls can fail on ungoverned supporting infrastructure.

Reflection prompt: If the control were disabled for two weeks, which indicators, owners, records and customer commitments would reveal the failure?

\Needspace{6\baselineskip}
\subsection{CC5.3 -- Policy deployment}
\begin{aicpabox}
\textbf{CC5.3}. COSO Principle 12: The entity deploys control activities through policies that establish what is expected and in procedures that put policies into action.\par\sourcepage{33}
\end{aicpabox}
Engineering interpretation (example): Convert approved policy into assigned, timed operating procedures and exceptions governance.

Evidence reasoning: Policies; procedures; schedules; exception register. Use this only after confirming the control is part of the actual system.

Assessment exercise: Select policies and trace to accountable operational executions. Identify the population, sample unit, expected outcome and exception-handling process.

Common misconception: Documents do not establish that activities occurred when required.

Reflection prompt: If the control were disabled for two weeks, which indicators, owners, records and customer commitments would reveal the failure?

\section{Logical and Physical Access Controls}
\Needspace{6\baselineskip}
\subsection{CC6.1 -- Logical access architecture}
\begin{aicpabox}
\textbf{CC6.1}. The entity implements logical access security software, infrastructure, and architectures over protected information assets to protect them from security events to meet the entity's objectives.\par\sourcepage{34}
\end{aicpabox}
Engineering interpretation (example): Federate identities; protect service identities; enforce MFA, RBAC and secured identity architecture.

Evidence reasoning: IAM configuration; architecture; identity inventory. Use this only after confirming the control is part of the actual system.

Assessment exercise: Inspect identities and authentication design across in-scope assets. Identify the population, sample unit, expected outcome and exception-handling process.

Common misconception: Service accounts and machine identities are often overlooked.

Reflection prompt: If the control were disabled for two weeks, which indicators, owners, records and customer commitments would reveal the failure?

\Needspace{6\baselineskip}
\subsection{CC6.2 -- Provisioning and deprovisioning}
\begin{aicpabox}
\textbf{CC6.2}. Prior to issuing system credentials and granting system access, the entity registers and authorizes new internal and external users whose access is administered by the entity. For those users whose access is administered by the entity, user system credentials are removed when user access is no longer authorized.\par\sourcepage{36}
\end{aicpabox}
Engineering interpretation (example): Bind user creation, unique identities and timely removal to approved joiner-mover-leaver workflow.

Evidence reasoning: Access requests; approval log; HR termination feed. Use this only after confirming the control is part of the actual system.

Assessment exercise: Sample additions and leavers from source population to access outcomes. Identify the population, sample unit, expected outcome and exception-handling process.

Common misconception: Ad hoc accounts that bypass central provisioning break the control.

Reflection prompt: If the control were disabled for two weeks, which indicators, owners, records and customer commitments would reveal the failure?

\Needspace{6\baselineskip}
\subsection{CC6.3 -- Least privilege review}
\begin{aicpabox}
\textbf{CC6.3}. The entity authorizes, modifies, or removes access to data, software, functions, and other protected information assets based on roles, responsibilities, or the system design and changes, giving consideration to the concepts of least privilege and segregation of duties, to meet the entity's objectives.\par\sourcepage{36}
\end{aicpabox}
Engineering interpretation (example): Use least-privilege roles, approval for privilege elevation, periodic access recertification and revocation.

Evidence reasoning: Access review signoffs; role policy; permissions exports. Use this only after confirming the control is part of the actual system.

Assessment exercise: Trace sampled privileged and ordinary access to approved role requirements. Identify the population, sample unit, expected outcome and exception-handling process.

Common misconception: An access export alone is not an approved recertification.

Reflection prompt: If the control were disabled for two weeks, which indicators, owners, records and customer commitments would reveal the failure?

\Needspace{6\baselineskip}
\subsection{CC6.4 -- Physical access}
\begin{aicpabox}
\textbf{CC6.4}. The entity restricts physical access to facilities and protected information assets (for example, data center facilities, backup media storage, and other sensitive locations) to authorized personnel to meet the entity's objectives.\par\sourcepage{37}
\end{aicpabox}
Engineering interpretation (example): Restrict facilities and backup-media access with appropriate badges, logs and outsourced control coverage.

Evidence reasoning: Badge access review; facility records; subservice reports. Use this only after confirming the control is part of the actual system.

Assessment exercise: Inspect applicable facility access and cloud provider responsibility boundary. Identify the population, sample unit, expected outcome and exception-handling process.

Common misconception: Office badges may not cover data-center physical security.

Reflection prompt: If the control were disabled for two weeks, which indicators, owners, records and customer commitments would reveal the failure?

\Needspace{6\baselineskip}
\subsection{CC6.5 -- Secure asset disposal}
\begin{aicpabox}
\textbf{CC6.5}. The entity discontinues logical and physical protections over physical assets only after the ability to read or recover data and software from those assets has been diminished and is no longer required to meet the entity's objectives.\par\sourcepage{37}
\end{aicpabox}
Engineering interpretation (example): Erase, destroy or make unrecoverable data and software before disposition; retain chain of custody.

Evidence reasoning: Sanitization records; disposal certificates; asset ledger. Use this only after confirming the control is part of the actual system.

Assessment exercise: Sample decommissioned assets to verify destruction and approval. Identify the population, sample unit, expected outcome and exception-handling process.

Common misconception: Asset recycling without verifiable sanitization leaves residual exposure.

Reflection prompt: If the control were disabled for two weeks, which indicators, owners, records and customer commitments would reveal the failure?

\Needspace{6\baselineskip}
\subsection{CC6.6 -- Perimeter and external threat defenses}
\begin{aicpabox}
\textbf{CC6.6}. The entity implements logical access security measures to protect against threats from sources outside its system boundaries.\par\sourcepage{38}
\end{aicpabox}
Engineering interpretation (example): Protect exposed interfaces with restricted network paths, identity controls and hardening.

Evidence reasoning: Firewall rules; segmentation map; WAF/config reviews. Use this only after confirming the control is part of the actual system.

Assessment exercise: Select exposed services and test approved ingress and egress posture. Identify the population, sample unit, expected outcome and exception-handling process.

Common misconception: Perimeter safeguards do not eliminate cloud public resource exposure.

Reflection prompt: If the control were disabled for two weeks, which indicators, owners, records and customer commitments would reveal the failure?

\Needspace{6\baselineskip}
\subsection{CC6.7 -- Secure data movement}
\begin{aicpabox}
\textbf{CC6.7}. The entity restricts the transmission, movement, and removal of information to authorized internal and external users and processes, and protects it during transmission, movement, or removal to meet the entity's objectives.\par\sourcepage{38}
\end{aicpabox}
Engineering interpretation (example): Classify data and restrict export, transmission and removable media; enforce suitable encryption.

Evidence reasoning: DLP policy; TLS configurations; export approvals. Use this only after confirming the control is part of the actual system.

Assessment exercise: Sample sensitive data transfers for authorized recipients and protection. Identify the population, sample unit, expected outcome and exception-handling process.

Common misconception: Encrypted transmission does not prove authorized disclosure.

Reflection prompt: If the control were disabled for two weeks, which indicators, owners, records and customer commitments would reveal the failure?

\Needspace{6\baselineskip}
\subsection{CC6.8 -- Malicious software prevention}
\begin{aicpabox}
\textbf{CC6.8}. The entity implements controls to prevent or detect and act upon the introduction of unauthorized or malicious software to meet the entity's objectives.\par\sourcepage{39}
\end{aicpabox}
Engineering interpretation (example): Prevent unapproved software and monitor malware on supported endpoints and workloads.

Evidence reasoning: EDR coverage; quarantine events; allowlist policy. Use this only after confirming the control is part of the actual system.

Assessment exercise: Reconcile endpoint population to protection status and exceptions. Identify the population, sample unit, expected outcome and exception-handling process.

Common misconception: An installed agent with disabled telemetry is not effective protection.

Reflection prompt: If the control were disabled for two weeks, which indicators, owners, records and customer commitments would reveal the failure?

\section{System Operations}
\Needspace{6\baselineskip}
\subsection{CC7.1 -- Vulnerability and drift detection}
\begin{aicpabox}
\textbf{CC7.1}. To meet its objectives, the entity uses detection and monitoring procedures to identify (1) changes to configurations that result in the introduction of new vulnerabilities, and (2) susceptibilities to newly discovered vulnerabilities.\par\sourcepage{40}
\end{aicpabox}
Engineering interpretation (example): Monitor vulnerability disclosures, baseline drift and asset exposure; prioritize and track remediation.

Evidence reasoning: Scan results; configuration diffs; risk acceptance tickets. Use this only after confirming the control is part of the actual system.

Assessment exercise: Sample vulnerable assets and follow discovery, ownership and remediation. Identify the population, sample unit, expected outcome and exception-handling process.

Common misconception: CodeQL alerts alone do not cover host, network and cloud exposures.

Reflection prompt: If the control were disabled for two weeks, which indicators, owners, records and customer commitments would reveal the failure?

\Needspace{6\baselineskip}
\subsection{CC7.2 -- Anomaly monitoring}
\begin{aicpabox}
\textbf{CC7.2}. The entity monitors system components and the operation of those components for anomalies that are indicative of malicious acts, natural disasters, and errors affecting the entity's ability to meet its objectives; anomalies are analyzed to determine whether they represent security events.\par\sourcepage{40}
\end{aicpabox}
Engineering interpretation (example): Collect and correlate infrastructure, application and identity signals with ownership and defined review.

Evidence reasoning: Detection rules; SIEM alerts; on-call histories. Use this only after confirming the control is part of the actual system.

Assessment exercise: Select detections and inspect triage and tuning evidence. Identify the population, sample unit, expected outcome and exception-handling process.

Common misconception: Telemetry collection without analysis is not sufficient monitoring.

Reflection prompt: If the control were disabled for two weeks, which indicators, owners, records and customer commitments would reveal the failure?

\Needspace{6\baselineskip}
\subsection{CC7.3 -- Security event evaluation}
\begin{aicpabox}
\textbf{CC7.3}. The entity evaluates security events to determine whether they could or have resulted in a failure of the entity to meet its objectives (security incidents) and, if so, takes actions to prevent or address such failures.\par\sourcepage{41}
\end{aicpabox}
Engineering interpretation (example): Classify events, assess impact, escalate incidents and preserve investigation records.

Evidence reasoning: Triage cases; severity rubric; incident timeline. Use this only after confirming the control is part of the actual system.

Assessment exercise: Sample events through classification and incident decision records. Identify the population, sample unit, expected outcome and exception-handling process.

Common misconception: Alert dismissal without documented investigation undermines event evaluation.

Reflection prompt: If the control were disabled for two weeks, which indicators, owners, records and customer commitments would reveal the failure?

\Needspace{6\baselineskip}
\subsection{CC7.4 -- Incident response}
\begin{aicpabox}
\textbf{CC7.4}. The entity responds to identified security incidents by executing a defined incident-response program to understand, contain, remediate, and communicate security incidents, as appropriate.\par\sourcepage{42}
\end{aicpabox}
Engineering interpretation (example): Practice detection, containment, eradication, coordination, communications and post-incident review.

Evidence reasoning: IR plan; tabletop results; incident cases. Use this only after confirming the control is part of the actual system.

Assessment exercise: Walk an incident through roles, approvals, communication and remediation. Identify the population, sample unit, expected outcome and exception-handling process.

Common misconception: A response playbook never exercised may be unproven.

Reflection prompt: If the control were disabled for two weeks, which indicators, owners, records and customer commitments would reveal the failure?

\Needspace{6\baselineskip}
\subsection{CC7.5 -- Security incident recovery}
\begin{aicpabox}
\textbf{CC7.5}. The entity identifies, develops, and implements activities to recover from identified security incidents.\par\sourcepage{44}
\end{aicpabox}
Engineering interpretation (example): Define restoration sequencing, integrity checks and post-incident validation and lessons learned.

Evidence reasoning: Recovery playbooks; restore tests; PIR action list. Use this only after confirming the control is part of the actual system.

Assessment exercise: Trace recovered services to tested restoration and business acceptance. Identify the population, sample unit, expected outcome and exception-handling process.

Common misconception: Restoring service without validating data integrity can prolong impact.

Reflection prompt: If the control were disabled for two weeks, which indicators, owners, records and customer commitments would reveal the failure?

\section{Change Management}
\Needspace{6\baselineskip}
\subsection{CC8.1 -- Controlled change management}
\begin{aicpabox}
\textbf{CC8.1}. The entity authorizes, designs, develops or acquires, configures, documents, tests, approves, and implements changes to infrastructure, data, software, and procedures to meet its objectives.\par\sourcepage{45}
\end{aicpabox}
Engineering interpretation (example): Use PR reviews, required checks, approvals, release provenance, rollback and emergency change handling.

Evidence reasoning: Pull requests; checks; releases; change tickets. Use this only after confirming the control is part of the actual system.

Assessment exercise: Sample normal and emergency production changes end-to-end. Identify the population, sample unit, expected outcome and exception-handling process.

Common misconception: Scan results after merge cannot establish a pre-merge control.

Reflection prompt: If the control were disabled for two weeks, which indicators, owners, records and customer commitments would reveal the failure?

\section{Risk Mitigation}
\Needspace{6\baselineskip}
\subsection{CC9.1 -- Business disruption mitigation}
\begin{aicpabox}
\textbf{CC9.1}. The entity identifies, selects, and develops risk mitigation activities for risks arising from potential business disruptions.\par\sourcepage{47}
\end{aicpabox}
Engineering interpretation (example): Define plausible disruption scenarios and proportionate mitigation including alternates and financial options.

Evidence reasoning: BIA; continuity response plan; mitigation decisions. Use this only after confirming the control is part of the actual system.

Assessment exercise: Trace business-impact findings to mitigations and assigned owners. Identify the population, sample unit, expected outcome and exception-handling process.

Common misconception: Disaster recovery alone does not address every disruption dependency.

Reflection prompt: If the control were disabled for two weeks, which indicators, owners, records and customer commitments would reveal the failure?

\Needspace{6\baselineskip}
\subsection{CC9.2 -- Third-party risk}
\begin{aicpabox}
\textbf{CC9.2}. The entity assesses and manages risks associated with vendors and business partners.\par\sourcepage{48}
\end{aicpabox}
Engineering interpretation (example): Inventory material vendors and assess security, contracts, shared responsibilities and ongoing performance.

Evidence reasoning: Vendor register; diligence; contracts; monitoring evidence. Use this only after confirming the control is part of the actual system.

Assessment exercise: Sample critical providers for due diligence and renewal reassessment. Identify the population, sample unit, expected outcome and exception-handling process.

Common misconception: A supplier's SOC 2 report does not automatically cover your obligations.

Reflection prompt: If the control were disabled for two weeks, which indicators, owners, records and customer commitments would reveal the failure?

\section{Additional Criteria for Availability}
\Needspace{6\baselineskip}
\subsection{A1.1 -- Capacity management}
\begin{aicpabox}
\textbf{A1.1}. The entity maintains, monitors, and evaluates current processing capacity and use of system components (infrastructure, data, and software) to manage capacity demand and to enable the implementation of additional capacity to help meet its objectives.\par\sourcepage{50}
\end{aicpabox}
Engineering interpretation (example): Track utilization and saturation, forecast load and establish thresholds for scaling and architecture changes.

Evidence reasoning: Capacity dashboards; forecast; scaling decisions. Use this only after confirming the control is part of the actual system.

Assessment exercise: Recompute capacity headroom and inspect change actions on breaches. Identify the population, sample unit, expected outcome and exception-handling process.

Common misconception: Availability promises need system-specific capacity assumptions.

Reflection prompt: If the control were disabled for two weeks, which indicators, owners, records and customer commitments would reveal the failure?

\Needspace{6\baselineskip}
\subsection{A1.2 -- Environmental and backup protections}
\begin{aicpabox}
\textbf{A1.2}. The entity authorizes, designs, develops or acquires, implements, operates, approves, maintains, and monitors environmental protections, software, data backup processes, and recovery infrastructure to meet its objectives.\par\sourcepage{50}
\end{aicpabox}
Engineering interpretation (example): Implement backup, power/environmental protection, recovery infrastructure, restoration plans and monitoring.

Evidence reasoning: Backup jobs; retention settings; facility/provider evidence. Use this only after confirming the control is part of the actual system.

Assessment exercise: Sample backups for success, protected storage and recovery dependencies. Identify the population, sample unit, expected outcome and exception-handling process.

Common misconception: Green backup jobs do not prove recoverability.

Reflection prompt: If the control were disabled for two weeks, which indicators, owners, records and customer commitments would reveal the failure?

\Needspace{6\baselineskip}
\subsection{A1.3 -- Recovery testing}
\begin{aicpabox}
\textbf{A1.3}. The entity tests recovery plan procedures supporting system recovery to meet its objectives.\par\sourcepage{52}
\end{aicpabox}
Engineering interpretation (example): Periodically exercise restoration procedures, including integrity checks, time objectives and follow-up.

Evidence reasoning: DR exercise plan; restore logs; issue remediation. Use this only after confirming the control is part of the actual system.

Assessment exercise: Reperform selected recovery steps and compare outcome with defined objectives. Identify the population, sample unit, expected outcome and exception-handling process.

Common misconception: Test scripts without execution results are not recovery evidence.

Reflection prompt: If the control were disabled for two weeks, which indicators, owners, records and customer commitments would reveal the failure?

\section{Additional Criteria for Confidentiality}
\Needspace{6\baselineskip}
\subsection{C1.1 -- Confidential information life cycle}
\begin{aicpabox}
\textbf{C1.1}. The entity identifies and maintains confidential information to meet the entity's objectives related to confidentiality.\par\sourcepage{52}
\end{aicpabox}
Engineering interpretation (example): Identify confidential datasets, restrict permitted use, retain appropriately and review protection needs.

Evidence reasoning: Classification inventory; data access matrix; retention map. Use this only after confirming the control is part of the actual system.

Assessment exercise: Sample data classes from creation through storage and access controls. Identify the population, sample unit, expected outcome and exception-handling process.

Common misconception: Confidentiality is broader than personally identifiable information.

Reflection prompt: If the control were disabled for two weeks, which indicators, owners, records and customer commitments would reveal the failure?

\Needspace{6\baselineskip}
\subsection{C1.2 -- Confidential disposal}
\begin{aicpabox}
\textbf{C1.2}. The entity disposes of confidential information to meet the entity's objectives related to confidentiality.\par\sourcepage{52}
\end{aicpabox}
Engineering interpretation (example): Trigger approved deletion, media destruction and vendor disposition at the end of retention.

Evidence reasoning: Retention schedule; deletion jobs; vendor certificates. Use this only after confirming the control is part of the actual system.

Assessment exercise: Sample expired records and verify timely irreversible disposition. Identify the population, sample unit, expected outcome and exception-handling process.

Common misconception: Deleting an application row may leave recoverable replicas and backups.

Reflection prompt: If the control were disabled for two weeks, which indicators, owners, records and customer commitments would reveal the failure?

\section{Additional Criteria for Processing Integrity}
\Needspace{6\baselineskip}
\subsection{PI1.1 -- Processing specifications}
\begin{aicpabox}
\textbf{PI1.1}. The entity obtains or generates, uses, and communicates relevant, quality information regarding the objectives related to processing, including definitions of data processed and product and service specifications, to support the use of products and services.\par\sourcepage{53}
\end{aicpabox}
Engineering interpretation (example): Define data contracts, schemas, field semantics, functional outcomes and user-facing specifications.

Evidence reasoning: Data dictionary; schema registry; acceptance criteria. Use this only after confirming the control is part of the actual system.

Assessment exercise: Trace specifications to implemented validations and published outputs. Identify the population, sample unit, expected outcome and exception-handling process.

Common misconception: Undocumented source units and populations cause apparent accuracy errors.

Reflection prompt: If the control were disabled for two weeks, which indicators, owners, records and customer commitments would reveal the failure?

\Needspace{6\baselineskip}
\subsection{PI1.2 -- Input completeness and accuracy}
\begin{aicpabox}
\textbf{PI1.2}. The entity implements policies and procedures over system inputs, including controls over completeness and accuracy, to result in products, services, and reporting to meet the entity's objectives.\par\sourcepage{54}
\end{aicpabox}
Engineering interpretation (example): Validate required fields, types, authorization, duplicates and completeness before processing.

Evidence reasoning: Validation rules; rejects; ingestion logs. Use this only after confirming the control is part of the actual system.

Assessment exercise: Inject invalid and duplicate data in controlled test and inspect handling. Identify the population, sample unit, expected outcome and exception-handling process.

Common misconception: Schema validation alone may not detect missing transactions.

Reflection prompt: If the control were disabled for two weeks, which indicators, owners, records and customer commitments would reveal the failure?

\Needspace{6\baselineskip}
\subsection{PI1.3 -- Processing correctness}
\begin{aicpabox}
\textbf{PI1.3}. The entity implements policies and procedures over system processing to result in products, services, and reporting to meet the entity's objectives.\par\sourcepage{54}
\end{aicpabox}
Engineering interpretation (example): Use idempotence, error handling, reconciliation and traceability for transformations.

Evidence reasoning: Job runs; test results; reconciliation outcomes. Use this only after confirming the control is part of the actual system.

Assessment exercise: Trace representative inputs through transformations and exception queues. Identify the population, sample unit, expected outcome and exception-handling process.

Common misconception: Successful batch exit codes do not prove complete processing.

Reflection prompt: If the control were disabled for two weeks, which indicators, owners, records and customer commitments would reveal the failure?

\Needspace{6\baselineskip}
\subsection{PI1.4 -- Output distribution}
\begin{aicpabox}
\textbf{PI1.4}. The entity implements policies and procedures to make available or deliver output completely, accurately, and timely in accordance with specifications to meet the entity's objectives.\par\sourcepage{55}
\end{aicpabox}
Engineering interpretation (example): Validate outputs, authorize recipients, preserve integrity and meet delivery expectations.

Evidence reasoning: Delivery manifests; approvals; integrity checks. Use this only after confirming the control is part of the actual system.

Assessment exercise: Trace approved outputs to recipients, timing and end-to-end completeness. Identify the population, sample unit, expected outcome and exception-handling process.

Common misconception: Accurate output sent to the wrong party fails the processing objective.

Reflection prompt: If the control were disabled for two weeks, which indicators, owners, records and customer commitments would reveal the failure?

\Needspace{6\baselineskip}
\subsection{PI1.5 -- Processing records retention}
\begin{aicpabox}
\textbf{PI1.5}. The entity implements policies and procedures to store inputs, items in processing, and outputs completely, accurately, and timely in accordance with system specifications to meet the entity's objectives.\par\sourcepage{56}
\end{aicpabox}
Engineering interpretation (example): Retain and protect input, intermediate and output records per system specification.

Evidence reasoning: Retention policies; transaction logs; storage tests. Use this only after confirming the control is part of the actual system.

Assessment exercise: Sample lifecycle records for accuracy, recoverability and defined retention. Identify the population, sample unit, expected outcome and exception-handling process.

Common misconception: Deleting intermediates without considering auditability can impair assurance.

Reflection prompt: If the control were disabled for two weeks, which indicators, owners, records and customer commitments would reveal the failure?

\section{Additional Criteria for Privacy: Notice}
\Needspace{6\baselineskip}
\subsection{P1.1 -- Privacy notice}
\begin{aicpabox}
\textbf{P1.1}. The entity provides notice to data subjects about its privacy practices to meet the entity's objectives related to privacy. The notice is updated and communicated to data subjects in a timely manner for changes to the entity's privacy practices, including changes in the use of personal information, to meet the entity's objectives related to privacy.\par\sourcepage{57}
\end{aicpabox}
Engineering interpretation (example): Maintain accurate notices mapped to actual data processing and communicate material updates.

Evidence reasoning: Privacy notice versions; publication history; data map. Use this only after confirming the control is part of the actual system.

Assessment exercise: Compare sampled collection paths with notice and update history. Identify the population, sample unit, expected outcome and exception-handling process.

Common misconception: A generic notice can miss actual third-party collection.

Reflection prompt: If the control were disabled for two weeks, which indicators, owners, records and customer commitments would reveal the failure?

\section{Additional Criteria for Privacy: Choice and Consent}
\Needspace{6\baselineskip}
\subsection{P2.1 -- Choice and consent}
\begin{aicpabox}
\textbf{P2.1}. The entity communicates choices available regarding the collection, use, retention, disclosure, and disposal of personal information to the data subjects and the consequences, if any, of each choice. Explicit consent for the collection, use, retention, disclosure, and disposal of personal information is obtained from data subjects or other authorized persons, if required. Such consent is obtained only for the intended purpose of the information to meet the entity's objectives related to privacy. The entity's basis for determining implicit consent for the collection, use, retention, disclosure, and disposal of personal information is documented.\par\sourcepage{58}
\end{aicpabox}
Engineering interpretation (example): Capture consent or lawful choice decisions with purpose, version, timestamp and withdrawal workflow.

Evidence reasoning: Consent ledger; UX screens; withdrawal logs. Use this only after confirming the control is part of the actual system.

Assessment exercise: Sample preferences to downstream enforcement and change history. Identify the population, sample unit, expected outcome and exception-handling process.

Common misconception: A checked box does not prove choice persists across integrations.

Reflection prompt: If the control were disabled for two weeks, which indicators, owners, records and customer commitments would reveal the failure?

\section{Additional Criteria for Privacy: Collection}
\Needspace{6\baselineskip}
\subsection{P3.1 -- Purpose-limited collection}
\begin{aicpabox}
\textbf{P3.1}. Personal information is collected consistent with the entity's objectives related to privacy.\par\sourcepage{59}
\end{aicpabox}
Engineering interpretation (example): Map personal information fields to authorized purposes and minimize collection by default.

Evidence reasoning: Data inventory; form design; API schemas. Use this only after confirming the control is part of the actual system.

Assessment exercise: Compare collected fields with documented purpose and approved scope. Identify the population, sample unit, expected outcome and exception-handling process.

Common misconception: Telemetry may collect extra identifiers outside primary forms.

Reflection prompt: If the control were disabled for two weeks, which indicators, owners, records and customer commitments would reveal the failure?

\Needspace{6\baselineskip}
\subsection{P3.2 -- Explicit consent collection}
\begin{aicpabox}
\textbf{P3.2}. For information requiring explicit consent, the entity communicates the need for such consent as well as the consequences of a failure to provide consent for the request for personal information and obtains the consent prior to the collection of the information to meet the entity's objectives related to privacy.\par\sourcepage{60}
\end{aicpabox}
Engineering interpretation (example): Where explicit consent is required, communicate consequences and preserve affirmative evidence.

Evidence reasoning: Consent artifacts; notices; audit logs. Use this only after confirming the control is part of the actual system.

Assessment exercise: Sample transactions requiring explicit consent for recorded decision. Identify the population, sample unit, expected outcome and exception-handling process.

Common misconception: Do not assert every data processing purpose requires consent.

Reflection prompt: If the control were disabled for two weeks, which indicators, owners, records and customer commitments would reveal the failure?

\section{Additional Criteria for Privacy: Use, Retention, and Disposal}
\Needspace{6\baselineskip}
\subsection{P4.1 -- Purpose limitation}
\begin{aicpabox}
\textbf{P4.1}. The entity limits the use of personal information to the purposes identified in the entity's objectives related to privacy.\par\sourcepage{60}
\end{aicpabox}
Engineering interpretation (example): Implement use restrictions, lineage and secondary-use approvals for personal information.

Evidence reasoning: Purpose map; access policy; data-use approvals. Use this only after confirming the control is part of the actual system.

Assessment exercise: Trace secondary uses to approved purpose and technical enforcement. Identify the population, sample unit, expected outcome and exception-handling process.

Common misconception: Reusing data for analytics may exceed the stated objective.

Reflection prompt: If the control were disabled for two weeks, which indicators, owners, records and customer commitments would reveal the failure?

\Needspace{6\baselineskip}
\subsection{P4.2 -- Personal-information retention}
\begin{aicpabox}
\textbf{P4.2}. The entity retains personal information consistent with the entity's objectives related to privacy.\par\sourcepage{60}
\end{aicpabox}
Engineering interpretation (example): Define retention clocks, lawful holds, processing-specific schedules and automated review.

Evidence reasoning: Retention schedule; legal holds; purge reports. Use this only after confirming the control is part of the actual system.

Assessment exercise: Sample aged records against approved retention and legal holds. Identify the population, sample unit, expected outcome and exception-handling process.

Common misconception: Infinite default retention undermines stated commitments.

Reflection prompt: If the control were disabled for two weeks, which indicators, owners, records and customer commitments would reveal the failure?

\Needspace{6\baselineskip}
\subsection{P4.3 -- Personal-information disposal}
\begin{aicpabox}
\textbf{P4.3}. The entity securely disposes of personal information to meet the entity's objectives related to privacy.\par\sourcepage{61}
\end{aicpabox}
Engineering interpretation (example): Implement verifiable deletion across primary stores, replicas and processor handoffs.

Evidence reasoning: Erasure tickets; deletion evidence; vendor confirmations. Use this only after confirming the control is part of the actual system.

Assessment exercise: Trace selected deletion requests through complete affected systems. Identify the population, sample unit, expected outcome and exception-handling process.

Common misconception: Backups and third parties often complicate deletion completion.

Reflection prompt: If the control were disabled for two weeks, which indicators, owners, records and customer commitments would reveal the failure?

\section{Additional Criteria for Privacy: Access}
\Needspace{6\baselineskip}
\subsection{P5.1 -- Data subject access}
\begin{aicpabox}
\textbf{P5.1}. The entity grants identified and authenticated data subjects the ability to access their stored personal information for review and, upon request, provides physical or electronic copies of that information to data subjects to meet the entity's objectives related to privacy. If access is denied, data subjects are informed of the denial and reason for such denial, as required, to meet the entity's objectives related to privacy.\par\sourcepage{61}
\end{aicpabox}
Engineering interpretation (example): Authenticate applicable requests, locate personal information and respond within documented parameters.

Evidence reasoning: DSAR register; identity check; response record. Use this only after confirming the control is part of the actual system.

Assessment exercise: Sample access requests for verification, completeness and response. Identify the population, sample unit, expected outcome and exception-handling process.

Common misconception: Overcollection during identity verification introduces new privacy risk.

Reflection prompt: If the control were disabled for two weeks, which indicators, owners, records and customer commitments would reveal the failure?

\Needspace{6\baselineskip}
\subsection{P5.2 -- Correction and amendment}
\begin{aicpabox}
\textbf{P5.2}. The entity corrects, amends, or appends personal information based on information provided by data subjects and communicates such information to third parties, as committed or required, to meet the entity's objectives related to privacy. If a request for correction is denied, data subjects are informed of the denial and reason for such denial to meet the entity's objectives related to privacy.\par\sourcepage{62}
\end{aicpabox}
Engineering interpretation (example): Support correction requests, document decisions and propagate changes as applicable.

Evidence reasoning: Correction tickets; updated records; vendor notices. Use this only after confirming the control is part of the actual system.

Assessment exercise: Trace completed and denied corrections with reason and propagation. Identify the population, sample unit, expected outcome and exception-handling process.

Common misconception: A frontend correction may leave stale downstream copies.

Reflection prompt: If the control were disabled for two weeks, which indicators, owners, records and customer commitments would reveal the failure?

\section{Additional Criteria for Privacy: Disclosure and Notification}
\Needspace{6\baselineskip}
\subsection{P6.1 -- Authorized disclosure}
\begin{aicpabox}
\textbf{P6.1}. The entity discloses personal information to third parties with the explicit consent of data subjects and such consent is obtained prior to disclosure to meet the entity's objectives related to privacy.\par\sourcepage{63}
\end{aicpabox}
Engineering interpretation (example): Gate outbound personal information sharing by lawful authorization and documented purpose.

Evidence reasoning: Data-sharing inventory; agreements; consent reference. Use this only after confirming the control is part of the actual system.

Assessment exercise: Sample disclosures for authorized basis before transfer. Identify the population, sample unit, expected outcome and exception-handling process.

Common misconception: Consent is criterion wording; local legal and engagement context must be assessed.

Reflection prompt: If the control were disabled for two weeks, which indicators, owners, records and customer commitments would reveal the failure?

\Needspace{6\baselineskip}
\subsection{P6.2 -- Disclosure accounting}
\begin{aicpabox}
\textbf{P6.2}. The entity creates and retains a complete, accurate, and timely record of authorized disclosures of personal information to meet the entity's objectives related to privacy.\par\sourcepage{63}
\end{aicpabox}
Engineering interpretation (example): Log approved personal information disclosures with recipient, data scope and date.

Evidence reasoning: Disclosure log; export records; reconciliation. Use this only after confirming the control is part of the actual system.

Assessment exercise: Reconcile sampled data exports to complete authorization ledger. Identify the population, sample unit, expected outcome and exception-handling process.

Common misconception: Manual spreadsheets often miss API-based transfers.

Reflection prompt: If the control were disabled for two weeks, which indicators, owners, records and customer commitments would reveal the failure?

\Needspace{6\baselineskip}
\subsection{P6.3 -- Unauthorized disclosure recording}
\begin{aicpabox}
\textbf{P6.3}. The entity creates and retains a complete, accurate, and timely record of detected or reported unauthorized disclosures (including breaches) of personal information to meet the entity's objectives related to privacy.\par\sourcepage{64}
\end{aicpabox}
Engineering interpretation (example): Record suspected and confirmed unauthorized disclosures with details and disposition.

Evidence reasoning: Incident register; disclosure timeline; investigation records. Use this only after confirming the control is part of the actual system.

Assessment exercise: Trace sampled events to documented records and closure rationale. Identify the population, sample unit, expected outcome and exception-handling process.

Common misconception: A detection event without privacy evaluation may be overlooked.

Reflection prompt: If the control were disabled for two weeks, which indicators, owners, records and customer commitments would reveal the failure?

\Needspace{6\baselineskip}
\subsection{P6.4 -- Privacy vendor commitments}
\begin{aicpabox}
\textbf{P6.4}. The entity obtains privacy commitments from vendors and other third parties who have access to personal information to meet the entity's objectives related to privacy. The entity assesses those parties' compliance on a periodic and as-needed basis and takes corrective action, if necessary.\par\sourcepage{64}
\end{aicpabox}
Engineering interpretation (example): Establish vendor privacy commitments, assess adherence and remediate misuse.

Evidence reasoning: DPA/contract clauses; vendor assessments; remediation log. Use this only after confirming the control is part of the actual system.

Assessment exercise: Sample processors for commitment, review and corrective action. Identify the population, sample unit, expected outcome and exception-handling process.

Common misconception: A confidentiality clause alone may omit required privacy obligations.

Reflection prompt: If the control were disabled for two weeks, which indicators, owners, records and customer commitments would reveal the failure?

\Needspace{6\baselineskip}
\subsection{P6.5 -- Vendor incident notifications}
\begin{aicpabox}
\textbf{P6.5}. The entity obtains commitments from vendors and other third parties with access to personal information to notify the entity in the event of actual or suspected unauthorized disclosures of personal information. Such notifications are reported to appropriate personnel and acted on in accordance with established incident-response procedures to meet the entity's objectives related to privacy.\par\sourcepage{65}
\end{aicpabox}
Engineering interpretation (example): Obtain escalation commitments and integrate third-party notifications into IR procedures.

Evidence reasoning: Notification terms; escalation contacts; vendor tests. Use this only after confirming the control is part of the actual system.

Assessment exercise: Examine sampled vendor incidents and reporting timelines. Identify the population, sample unit, expected outcome and exception-handling process.

Common misconception: Untested contractual notice paths can delay regulatory actions.

Reflection prompt: If the control were disabled for two weeks, which indicators, owners, records and customer commitments would reveal the failure?

\Needspace{6\baselineskip}
\subsection{P6.6 -- Privacy incident notification}
\begin{aicpabox}
\textbf{P6.6}. The entity provides notification of breaches and incidents to affected data subjects, regulators, and others to meet the entity's objectives related to privacy.\par\sourcepage{65}
\end{aicpabox}
Engineering interpretation (example): Assess recipients and deadlines from applicable commitments and laws; document notification decisions.

Evidence reasoning: Breach decision tree; counsel review; notices. Use this only after confirming the control is part of the actual system.

Assessment exercise: Trace incidents to legal analysis and timely notification or rationale. Identify the population, sample unit, expected outcome and exception-handling process.

Common misconception: Do not assume every security event is a notifiable privacy breach.

Reflection prompt: If the control were disabled for two weeks, which indicators, owners, records and customer commitments would reveal the failure?

\Needspace{6\baselineskip}
\subsection{P6.7 -- Accounting and disclosure to subject}
\begin{aicpabox}
\textbf{P6.7}. The entity provides data subjects with an accounting of the personal information held and disclosure of the data subjects' personal information, upon the data subjects' request, to meet the entity's objectives related to privacy.\par\sourcepage{65}
\end{aicpabox}
Engineering interpretation (example): Provide appropriate accounting of held and disclosed personal information using verified records.

Evidence reasoning: DSAR package; inventory; disclosure accounting. Use this only after confirming the control is part of the actual system.

Assessment exercise: Sample requests for completeness, identity verification and response. Identify the population, sample unit, expected outcome and exception-handling process.

Common misconception: An incomplete disclosure inventory makes accounting inaccurate.

Reflection prompt: If the control were disabled for two weeks, which indicators, owners, records and customer commitments would reveal the failure?

\section{Additional Criteria for Privacy: Quality}
\Needspace{6\baselineskip}
\subsection{P7.1 -- Personal data quality}
\begin{aicpabox}
\textbf{P7.1}. The entity collects and maintains accurate, up-to-date, complete, and relevant personal information to meet the entity's objectives related to privacy.\par\sourcepage{66}
\end{aicpabox}
Engineering interpretation (example): Define validity, currency and completeness rules with correction and synchronization.

Evidence reasoning: Data-quality dashboard; reconciliations; corrected records. Use this only after confirming the control is part of the actual system.

Assessment exercise: Sample key attributes for quality controls and exceptions. Identify the population, sample unit, expected outcome and exception-handling process.

Common misconception: Security controls do not necessarily establish data accuracy.

Reflection prompt: If the control were disabled for two weeks, which indicators, owners, records and customer commitments would reveal the failure?

\section{Additional Criteria for Privacy: Monitoring and Enforcement}
\Needspace{6\baselineskip}
\subsection{P8.1 -- Privacy inquiry and dispute handling}
\begin{aicpabox}
\textbf{P8.1}. The entity implements a process for receiving, addressing, resolving, and communicating the resolution of inquiries, complaints, and disputes from data subjects and others and periodically monitors compliance to meet the entity's objectives related to privacy. Corrections and other necessary actions related to identified deficiencies are made or taken in a timely manner.\par\sourcepage{66}
\end{aicpabox}
Engineering interpretation (example): Offer intake, investigation, resolution, appeals and trend monitoring for privacy inquiries.

Evidence reasoning: Case register; SLA results; corrective actions. Use this only after confirming the control is part of the actual system.

Assessment exercise: Select disputes and examine receipt, handling, decision and follow-up. Identify the population, sample unit, expected outcome and exception-handling process.

Common misconception: Closing complaints without root-cause review hides systemic failures.

Reflection prompt: If the control were disabled for two weeks, which indicators, owners, records and customer commitments would reveal the failure?

\chapter{Cross-category application exercises}
\begin{enumerate}
\item An application provider promises confidentiality but has no documented classification, retention periods or vendor data map. Relate the gap to common criteria as well as the C series and decide what must be tested.
\item A SaaS platform has a fully automated deployment pipeline, but its CodeQL checks run only after merging to the default branch. Evaluate the distinction between a detection control and a pre-merge change authorization gate.
\item An organization records consent on its frontend, but its analytics pipeline does not ingest withdrawal updates. Trace the relevant privacy criteria and the evidence needed to demonstrate enforcement.
\item A backup service reports 100\% successful nightly jobs. Evaluate whether A1.2, A1.3 and recovery-related common criteria can be tested solely from that metric.
\item A service organization is a processor for customer data and a controller for its employee data. Evaluate which privacy points of focus are likely relevant in each processing context.
\end{enumerate}

\chapter{Self-check and model responses}
\subsection{Q1. Does a point of focus create a stand-alone compulsory control?}
Model response: No. Paragraph .07 explicitly states that addressing each point of focus need not be assessed; the relevant criterion and the entity-specific circumstances control the analysis.

\subsection{Q2. Does SOC 2 Type 1 express an opinion on operating effectiveness?}
Model response: No. The source contrasts a Type 1 report with Type 2: Type 2 includes an opinion on operating effectiveness and details of control testing during the specified period.

\subsection{Q3. Is privacy interchangeable with confidentiality?}
Model response: No. Privacy is about personal information and its collection, use, retention, disclosure and disposal; confidentiality can extend to nonpersonal protected information such as trade secrets.

\subsection{Q4. Does a scanner alert count as control completion?}
Model response: Not by itself. The organization must demonstrate that its defined detection and remediation process acted on relevant results, with evidence and risk decisions.

\subsection{Q5. What does [P][C] mean in the privacy-focused headings?}
Model response: Typical relevance to data processors and controllers. The annotations help evaluation but do not replace examining the actual relationship and engagement scope.


\chapter*{Source and editorial protocol}
\addcontentsline{toc}{chapter}{Source and editorial protocol}
This independent document is grounded in the user-supplied \emph{AICPA TSP Section 100: 2017 Trust Services Criteria for Security, Availability, Processing Integrity, Confidentiality, and Privacy (With Revised Points of Focus---2022)}, 75-page PDF, especially paragraphs .01--.29, the criterion table on pp.~14--67, and the glossary. Criterion wording is quoted when labeled as such. Selected point-of-focus labels are aids to interpretation and may be incomplete. Engineering recommendations, role assignments, tools, tests, dates, scorecards, and example frequencies are author-developed examples and are not AICPA requirements. The 2022 revisions changed points of focus, not the 2017 criteria (source p.~3). The source states that an assessment of each point of focus is not required (p.~5). This document does not constitute a CPA attestation or legal advice.
\end{document}
